% !TEX root = ../main.tex
A Szakdolgozat 1 során áttekintettük a mobiltelefon-hálózatok felépítését, a CDR
adatok szerkezetét és adatvédelmi kérdéseit. Részletes szakirodalmi áttekintést
nyújtottunk a CDR alapú humán mobilitás kutatásokról: bemutattuk, hogy a
távolságok eloszlása vágott hatványfüggvényt (truncated power-law) követ, és hogy
a mozgástér sugara jellemzi az egyének által bejárt területet \cite{gonzalez2008understanding}.
Elemeztük azokat az alkalmazási területeket, ahol a CDR hasznos lehet, mint
például a közlekedéstervezés, a smart city koncepciók \cite{batty2012smart}, a
szocioökonómiai elemzések \cite{pinter2021evaluating, rhoads2020measuring} és az
eseménydetektálás \cite{traag2011social}. Részletesen áttekintettük a mobilitási
metrikák számítási módszereit, különös tekintettel a mozgástér sugarára és az
entrópia-alapú mutatókra, és összefoglaltuk a fontosabb metrikákat egy táblázatban.

Bemutattuk továbbá a kutatási célú CDR adatkészletek általános struktúráját,
kitértünk az adatminőségre és az etikai megfontolásokra a szakirodalom alapján
\cite{pinter2022awakening, zhang2019connected}. Végül részletes rendszertervet dolgoztunk
ki a CDR adatok feldolgozására: megfogalmaztuk a funkcionális és nem
funkcionális követelményeket, vázoltuk a magas szintű architektúrát, kidolgoztuk
a csillagsémás adatbázis-tervet, az ETL folyamatot és a tervezett technológiák listáját.

\subsection{Kapcsolat a Szakdolgozat 2-vel}
A következő fázisban (Szakdolgozat 2) a tervezett rendszer gyakorlati megvalósítása
történik meg. Ennek főbb lépései:
\begin{itemize}
	\item A CDR adatok tényleges betöltése a staging és DWH környezetbe; az ETL
		folyamatok implementálása Python és SQL használatával.

	\item A mobilitási metrikák számítását végző modulok fejlesztése, beleértve a
		mozgástér sugara, a térbeli és időbeli entrópia, a home–work távolság és
		további mutatók implementálását
		\cite{pinter2021evaluating, pinter2022awakening}.

	\item A \texttt{fact\_mobility} tábla feltöltése, majd a mutatók exportja QGIS
		kompatibilis formátumba; tematikus térképek, hőtérképek és interaktív dashboardok
		készítése.

	\item Esettanulmányok készítése: például a város különböző kerületeinek
		összehasonlítása a mozgástér sugara vagy az entrópia alapján, ingázási hálózatok
		vizualizálása, valamint eseménydetektálás implementálása
		\cite{pinter2021evaluating, traag2011social}.

	\item A rendszer teljesítményének és pontosságának értékelése; skálázhatósági
		tesztek végrehajtása.
\end{itemize}
E munka eredményeként létrejön egy működő CDR analitikai rendszer, amely képes
valós városi adatok feldolgozására, a mobilitási mintázatok feltárására és térképi
megjelenítésére.


\subsection{A megvalósítás kiemelt kutatási iránya: Adatvezérelt szolgáltatás-tervezés}

A Szakdolgozat 2 implementációs fázisában a rendszer fejlesztését egy konkrét, piaci és városüzemeltetési szempontból is releváns koncepció köré építem fel: \textbf{„Városi ingázás-optimalizálás és dinamikus szolgáltatás-tervezés”}.

Jelenleg a várostervezés gyakran statikus (pl. népszámlálási) adatokra támaszkodik, amelyek nem követik le a valós idejű lakossági mozgásokat. A kutatásom célja, hogy a CDR adatok dinamikus természetét kihasználva egy olyan döntéstámogató réteget hozzak létre, amely feltárja a város rejtett infrastrukturális hiányosságait.

Ezt a célt a már definiált metrikák újszerű összekapcsolásával valósítom meg:

\begin{itemize}
    \item \textbf{Összefüggés-vizsgálat:} A \textit{Home-Work távolság} (mint kötött pályás mozgás) és a \textit{Radius of Gyration} ($r_g$) (mint teljes mozgástér) korrelációjának elemzése.
    \item \textbf{„Mobilitási kényszer” azonosítása:} Olyan városi zónák detektálása, ahol a lakók a feltételezettnél nagyobb mobilitási sugarat ($r_g$) kénytelenek bejárni, ami a helyi szolgáltatások (munkahelyek, kereskedelem) hiányára utal.
\end{itemize}

A megvalósított rendszer két konkrét területen nyújt majd üzleti és társadalmi értéket:

\begin{enumerate}
    \item \textbf{Közlekedésszervezési optimalizálás:} A statikus menetrendek helyett a valós csúcsidőszaki terhelések és a tényleges (nem lakcímkártya alapú) ingázási útvonalak feltárása, támogatva a járatsűrűség hatékonyabb tervezését.
    \item \textbf{Üzleti telephely-választás (Site Selection):} A rendszer képes lesz kijelölni azokat az „ellátatlan” zónákat, ahol magas a nappali népsűrűség, de aránytalanul nagy az utazási kényszer. Ez értékes inputot jelent kereskedelmi egységek vagy szolgáltatók számára az ideális lokáció kiválasztásához.
\end{enumerate}

Ez a fókusz biztosítja, hogy az elkészülő informatikai megoldás nem csupán technológiai demonstráció lesz, hanem valós városi problémákra (torlódás, szolgáltatáshiány) ad adatvezérelt válaszokat.
